\subsection{Faithfully flat rings and finiteness conditions}

PROPOSITION 11. Let B be a ring and A a subring of B such that B is a faithfully flat right A-module. For a left A-module F to be finitely generated (resp. finitely presented), it is necessary and sufficient that the B-module \(B \otimes_{A} F\) be finitely generated (resp. finitely presented).

\begin{enumerate}
\def\labelenumi{(\arabic{enumi})}
\item
  Without any hypothesis on B, clearly, if F is a finitely generated left A-module, \(B \otimes_{A} F\) is a finitely generated left B-module. Conversely, if \(B \otimes_{A} F\) is a finitely generated B-module, it is generated by a finite number of elements of the form \(1 \otimes x_{i}\) with \(x_{i} \in F\) ; if M is a sub-A-module of F generated by the \(x_{i}\) and j the canonical injection \(M \to F\) , \(1_{B} \otimes j: B \otimes_{A} M \to B \otimes_{A} F\) is a surjective homomorphism, hence j is surjective (no. 1, Proposition 2), which proves that F is finitely generated.
\item
  If \(\mathbf{F}\) admits a finite presentation, so does \(\mathbf{B} \otimes_{\mathbf{A}} \mathbf{F}\) without any hypothesis on B (§ 2, no. 8). It remains to prove that, if \(B \otimes_{A} F\) admits a finite presentation, so does F. We already know from (1) that F is finitely generated, hence there exists a surjective homomorphism u: L → F, where L is a finitely generated free A-module. Let R be the kernel of u, so that \(B \otimes_{A} R\) is identified with the kernel of the surjective homomorphism \(1_{B} \otimes u: B \otimes_{A} L \to B \otimes_{A} F\) (§ 2, no. 3, Remark 2). As \(B \otimes_{A} F\) admits a finite presentation by hypothesis, we conclude (§ 2, no. 8, Lemma 9) that \(B \otimes_{A} R\) is finitely generated; then it follows from (1) that R is a finitely generated A-module and consequently F admits a finite presentation.
\end{enumerate}

PROPOSITION 12. Let B be a ring and A a commutative subring of the centre of B such that B is a faithfully flat A-module. For an A-module F to be projective and finitely generated, it is necessary and sufficient that \(B \otimes_{A} F\) be a finitely generated projective left B-module.

The condition is obviously necessary without any hypothesis on A or B (Algebra, Chapter II, § 5, no. 1, Corollary to Proposition 4); we prove that it is sufficient. If a finitely generated projective module admits a finite presentation (§ 2, no. 8, Lemma 8), the hypothesis implies that F admits a finite presentation by virtue of Proposition 11, hence, for every A-module M, there is a canonical isomorphism

\[
\omega : \mathbf {B} \otimes_ {\mathbf {A}} \operatorname{Hom} _ {\mathbf {A}} (\mathbf {F}, \mathbf {M}) \rightarrow \operatorname{Hom} _ {\mathbf {B}} (\mathbf {B} \otimes_ {\mathbf {A}} \mathbf {F}, \mathbf {B} \otimes_ {\mathbf {A}} \mathbf {M})
\]

(§ 2, no. 10, Proposition 11). Then let \(v: \mathbf{M} \to \mathbf{M}''\) be a surjective A-module homomorphism and consider the commutative diagram

\[
\begin{array}{c} \mathrm{B} \otimes_ {\mathrm{A}} \operatorname{Hom} _ {\mathrm{A}} (\mathrm{F}, \mathrm{M}) \xrightarrow {\omega} \operatorname{Hom} _ {\mathrm{B}} (\mathrm{B} \otimes_ {\mathrm{A}} \mathrm{F}, \mathrm{B} \otimes_ {\mathrm{A}} \mathrm{M}) \\ 1 _ {\mathrm{B}} \otimes \operatorname{Hom} (1 _ {\mathrm{F}}, v) \Bigg \downarrow \quad \Bigg \downarrow \operatorname{Hom} (1 _ {\mathrm{B}} \otimes \mathrm{F}, 1 _ {\mathrm{B}} \otimes v) \\ \mathrm{B} \otimes_ {\mathrm{A}} \operatorname{Hom} _ {\mathrm{A}} (\mathrm{F}, \mathrm{M} ^ {\prime \prime}) \xrightarrow [ \omega ]{\longrightarrow} \operatorname{Hom} _ {\mathrm{B}} (\mathrm{B} \otimes_ {\mathrm{A}} \mathrm{F}, \mathrm{B} \otimes_ {\mathrm{A}} \mathrm{M} ^ {\prime \prime}) \end{array}
\]

As \(1_{B} \otimes v\) is surjective and \(B \otimes_{A} F\) is assumed projective, \(\operatorname{Hom}(1_{B} \otimes F, 1_{B} \otimes v)\) is surjective (Algebra, Chapter II, § 2, no. 2, Proposition 4) and so then is \(1_{B} \otimes \operatorname{Hom}(1_{F}, v)\) . But as B is a faithfully flat A-module, \(\operatorname{Hom}(1_{F}, v)\) is itself surjective (no. 1, Proposition 2), hence F is a projective A-module (Algebra, Chapter II, § 2, no. 2, Proposition 4).

